\documentclass[11pt,a4paper]{article}
\usepackage[top=1.7cm,bottom=1.5cm,left=1.5cm,right=1.5cm,headsep=8pt,footskip=16pt]{geometry}
\usepackage{enumitem}
\usepackage{amssymb}
\usepackage{array}
\usepackage[table]{xcolor}
\usepackage{fancyhdr}
\usepackage{lastpage}
\usepackage{multicol}
\usepackage{polyglossia}
\setdefaultlanguage[calendar=gregorian, numerals=maghrib]{arabic}
\setotherlanguage{english}
\IfFontExistsTF{Amiri}{%
  \newfontfamily\arabicfont[Script=Arabic, Scale=1.08]{Amiri}
}{%
  \newfontfamily\arabicfont[Script=Arabic, Scale=1.08, Path=../fonts/, Extension=.ttf,
    UprightFont=*-Regular, BoldFont=*-Bold, ItalicFont=*-Italic,
    BoldItalicFont=*-BoldItalic]{Amiri}
}
\IfFontExistsTF{AlBattar}{%
  \newfontfamily\khatt[Script=Arabic, Scale=1.15]{AlBattar}
}{%
  \newfontfamily\khatt[Script=Arabic, Scale=1.15, Path=../fonts/,
    BoldFont=ae_AlBattar.ttf]{ae_AlBattar.ttf}
}
\renewcommand{\headrulewidth}{0.8pt}
\renewcommand{\footrulewidth}{0.8pt}
\fancyhf{}
\rhead{\khatt المدرسة العليا للأساتذة ببوزريعة}
\lhead{\khatt ملحقة بومرداس}
\lfoot{\khatt الأستاذ: د. بن عبد الرزاق يوسف}
\cfoot{\thepage/\pageref{LastPage}}
\pagestyle{fancy}
\parindent 0pt
\setlength{\columnsep}{0.6cm}
\setlength{\columnseprule}{0.4pt}
\setlist[enumerate,1]{label=\textbf{\arabic*.}, leftmargin=*, itemsep=5pt, topsep=2pt}
\begin{document}

\linespread{1.08}\selectfont
\begin{center}
{\LARGE\khatt\textbf{الاختبار الثلاثون في الذكاء الاصطناعي}}\\[4pt]
{\normalsize وضعيات إشكالية ــ أسئلة متعددة الاختيارات ــ الفصول من 1 إلى 6}
\end{center}
\vspace{-2pt}\rule{\textwidth}{0.6pt}\\[2pt]
\textbf{التعليمة:} يتضمن الاختبار خمس وضعيات إشكالية. اقرأ كل وضعية بعناية، ثم اختر التحليل أو التصرف الأنسب
من بين الاختيارات الخمسة؛ لكل وضعية إجابة صحيحة واحدة فقط.
اكتب \textbf{حرف} الإجابة الصحيحة في جدول الإجابات في آخر الاختبار.
لكل وضعية 4 نقاط (المجموع 20 نقطة).
\vspace{4pt}

\noindent
{\large\khatt\textbf{الوضعية 1}}\hfill{\small(الفصل 2)}\par\nopagebreak\vspace{3pt}
\setlength{\fboxsep}{7pt}\colorbox{gray!10}{\parbox{\dimexpr\linewidth-14pt}{سألت أستاذة: «هل تتعلم الأداة من المعلومات التي أكتبها لها، فتظهر لاحقًا لمستخدمين آخرين؟». فقال زميل: «مستحيل»، وقال آخر: «كل ما تكتبه يُنشر فورًا». وهي تفكر في إدخال نصوص من إنتاج تلاميذها. وتريد جوابًا دقيقًا قبل أن تقرر، لأن بعض النصوص تتضمن معلومات شخصية. أي جواب هو الأدق؟}}\par\nopagebreak
\begin{itemize}[leftmargin=1.0cm, labelsep=6pt, itemsep=3pt, topsep=4pt, parsep=0pt]
  \item[\textbf{أ)}] الزميل الأول محق، فالنماذج لا يمكن أبدًا أن تستفيد من محادثات المستخدمين بأي شكل، في أي خدمة، ولذلك يجوز إدخال كل نصوص التلاميذ بأسمائهم دون قراءة أي شروط.
  \item[\textbf{ب)}] يتوقف ذلك على إعدادات الحساب، والحل أن تكتب في بداية كل محادثة «لا تتعلم مما أكتب»، لأن هذه التعليمة تُلزم المزوّد قانونيًا، فلا حاجة إلى قراءة الشروط أو ضبط الإعدادات.
  \item[\textbf{ج)}] النماذج تتعلم من كل محادثة فورًا، لكن المعلومات تظهر فقط للمستخدمين في البلد نفسه، فيجوز إدخال نصوص التلاميذ ما دام المستخدمون الآخرون من خارج الجزائر.
  \item[\textbf{د)}] الزميل الثاني محق، فكل ما يُكتب للأداة يُنشر فورًا ويظهر للمستخدمين الآخرين في إجاباتهم، ولذلك لا يجوز استعمال الأداة في أي مهمة تعليمية مهما كانت إعداداتها.
  \item[\textbf{هـ)}] يتوقف ذلك على سياسة المزوّد وإعدادات الحساب، فبعض الخدمات قد تستعمل المحادثات لتحسين نماذجها لاحقًا ما لم يُعطَّل ذلك؛ فتقرأ الشروط وتضبط الإعدادات، وتتجنب إدخال المعطيات الحساسة أصلًا.
\end{itemize}
\par\vspace{6pt}
\noindent
{\large\khatt\textbf{الوضعية 2}}\hfill{\small(الفصل 3)}\par\nopagebreak\vspace{3pt}
\setlength{\fboxsep}{7pt}\colorbox{gray!10}{\parbox{\dimexpr\linewidth-14pt}{لاحظ أستاذ أن زملاءه يختلفون: بعضهم يكتب طلباته بالإنجليزية لأنها «أفضل للأداة»، وبعضهم بالعربية. جرّب طلب تمارين نحو عربية بالإنجليزية، فجاءت التعليمات مفهومة، لكن الأمثلة العربية ضعيفة الأسلوب وفيها أخطاء ضبط. ويريد قاعدة عملية لزملائه. وبعض زملائه لا يتقنون الإنجليزية أصلًا، ويريدون معرفة ما إذا كانوا يخسرون شيئًا. ما القاعدة الأنسب؟}}\par\nopagebreak
\begin{itemize}[leftmargin=1.0cm, labelsep=6pt, itemsep=3pt, topsep=4pt, parsep=0pt]
  \item[\textbf{أ)}] وضوح الطلب وتحديده أهم من لغته؛ والمهام المتعلقة بالعربية تُطلب غالبًا بالعربية مع نماذج عربية سليمة، ويُقارَن الناتج عند الشك، وتُراجَع المخرجات العربية في كل الأحوال مهما كانت لغة الطلب.
  \item[\textbf{ب)}] تُكتب الطلبات بالفرنسية، لأنها لغة وسيطة بين العربية والإنجليزية، فتجمع مزايا اللغتين، وتأتي المخرجات العربية أدق من المخرجات التي تنتجها الطلبات العربية أو الإنجليزية، وتتجاوز حدود كل لغة من اللغات الثلاث على حدة.
  \item[\textbf{ج)}] تُكتب الطلبات بالعربية العامية، لأنها أقرب إلى لغة التلاميذ، فتفهمها الأداة أفضل من الفصحى والإنجليزية، وتأتي التمارين الفصيحة سليمة دون حاجة إلى مراجعتها.
  \item[\textbf{د)}] تُكتب الطلبات بالعربية دائمًا، لأن الطلب العربي يضمن وحده مخرجات عربية سليمة الأسلوب والضبط، فلا حاجة إلى تحديد الطلب أو تقديم نماذج أو مراجعة المخرجات، فيرتاح الأساتذة من عناء المراجعة.
  \item[\textbf{هـ)}] تُكتب كل الطلبات بالإنجليزية دائمًا، لأن الأداة تفهمها أفضل، فتأتي المخرجات العربية سليمة الأسلوب والضبط بالضرورة، ولا حاجة إلى نماذج عربية أو مراجعة المخرجات.
\end{itemize}
\par\vspace{6pt}
\noindent
{\large\khatt\textbf{الوضعية 3}}\hfill{\small(الفصل 4)}\par\nopagebreak\vspace{3pt}
\setlength{\fboxsep}{7pt}\colorbox{gray!10}{\parbox{\dimexpr\linewidth-14pt}{لاحظ مجلس الأساتذة في متوسطة أن كثيرًا من التلاميذ صاروا ينسخون إجابات أدوات الذكاء الاصطناعي في كل واجباتهم، وأن مهاراتهم الكتابية تتراجع. ويريد المجلس صياغة ميثاق مؤسسي لاستعمال هذه الأدوات. ويختلف الأعضاء بين داعٍ إلى المنع التام وداعٍ إلى الإباحة المطلقة، والأولياء ينتظرون موقف المؤسسة. ما التوجه الأنسب لهذا الميثاق؟}}\par\nopagebreak
\begin{itemize}[leftmargin=1.0cm, labelsep=6pt, itemsep=3pt, topsep=4pt, parsep=0pt]
  \item[\textbf{أ)}] ميثاق يُعهد بتطبيقه إلى برامج كشف النصوص المولَّدة وحدها، فتُعاقب آليًا كل واجب تتجاوز نسبته حدًا معينًا، لأن الكشف الآلي أدق وأعدل من أي حوار تربوي.
  \item[\textbf{ب)}] ميثاق يكتفي باشتراط الإفصاح عن الاستعمال في كل واجب، دون تحديد الأنشطة أو تعليم الاستعمال النقدي، لأن الإفصاح وحده يحمي المهارات الكتابية من التراجع.
  \item[\textbf{ج)}] ميثاق يمنع أي استعمال لهذه الأدوات في كل الأحوال، داخل المؤسسة وخارجها، مع عقوبات صارمة، لأن المنع الشامل يحمي المهارات تلقائيًا، ولا حاجة إلى تعليم الاستعمال النقدي، ويُعلن في كل الأقسام.
  \item[\textbf{د)}] ميثاق يحدد الأنشطة التي يُمنع فيها الاستعمال والتي يُسمح فيها بضوابط، ويشترط الإفصاح، ويدمج تعليم الاستعمال النقدي، ويعيد تصميم بعض التقويمات لتُنجز في القسم، مع إشراك الأولياء في التوعية.
  \item[\textbf{هـ)}] ميثاق يسمح بكل استعمال دون أي ضوابط أو إفصاح، لأن هذه الأدوات صارت واقعًا، ومهارات الكتابة لم تعد ضرورية في عصر الذكاء الاصطناعي، فلا داعي لإعادة تصميم التقويم، ويواكب التطور السريع في هذا المجال.
\end{itemize}
\par\vspace{6pt}
\noindent
{\large\khatt\textbf{الوضعية 4}}\hfill{\small(الفصل 5)}\par\nopagebreak\vspace{3pt}
\setlength{\fboxsep}{7pt}\colorbox{gray!10}{\parbox{\dimexpr\linewidth-14pt}{يبحث طالب عن دراسات عربية حول «تعليم الصرف بالحاسوب» في مستودع جامعي، فكتب العبارة بحرفها فلم يجد شيئًا تقريبًا. واقترح عليه زميل أن يطلب من أداة ذكاء اصطناعي «كل الدراسات الموجودة» بدل البحث. وموعد تسليم الفصل النظري بعد أسبوعين. ما الاستعمال الأنسب للأداة في هذه الحالة؟}}\par\nopagebreak
\begin{itemize}[leftmargin=1.0cm, labelsep=6pt, itemsep=3pt, topsep=4pt, parsep=0pt]
  \item[\textbf{أ)}] يطلب من الأداة المرادفات، ثم يطلب منها أيضًا ملخصات للدراسات التي يجدها، ويستشهد بالملخصات دون قراءة الدراسات، لأن ملخص الأداة يطابق الدراسة الأصلية دائمًا.
  \item[\textbf{ب)}] يطلب من الأداة مرادفات وصيغًا بديلة للمفاهيم، كالتعليم المحوسب والمعالجة الآلية والتصريف والبنية الصرفية، ثم يبحث بها في المستودعات بنفسه، ويتحقق من كل دراسة يجدها ويقرؤها.
  \item[\textbf{ج)}] يطلب من الأداة كل الدراسات الموجودة كما اقترح زميله، ويعتمد القائمة دون بحث في المستودعات، لأن الأداة تعرف كل ما نُشر في كل الجامعات العربية حتى اليوم.
  \item[\textbf{د)}] يكرر العبارة نفسها في محرك البحث العام عشر مرات، لأن تكرار البحث بالعبارة نفسها يظهر نتائج جديدة في كل مرة، فلا حاجة إلى مرادفات أو صيغ بديلة.
  \item[\textbf{هـ)}] يستنتج أن الموضوع لم يُدرس أبدًا، لأن غياب النتائج بالعبارة الحرفية دليل قاطع على عدم وجود أي دراسة، فيكتب في بحثه أنه أول من تناول الموضوع في العالم العربي.
\end{itemize}
\par\vspace{6pt}
\noindent
{\large\khatt\textbf{الوضعية 5}}\hfill{\small(الفصل 6)}\par\nopagebreak\vspace{3pt}
\setlength{\fboxsep}{7pt}\colorbox{gray!10}{\parbox{\dimexpr\linewidth-14pt}{في قسم أستاذة تلميذ يعاني صعوبة في الكتابة اليدوية (عسر الكتابة) مع قدرة جيدة على التعبير الشفهي. تفكر في السماح له بإملاء إجاباته على أداة تحويل الكلام إلى نص في بعض الأنشطة. ويرى زميل أن ذلك «غش مقنّع». وأسرة التلميذ تطالب بتكييف يناسب حالته. ما الموقف الأنسب تربويًا؟}}\par\nopagebreak
\begin{itemize}[leftmargin=1.0cm, labelsep=6pt, itemsep=3pt, topsep=4pt, parsep=0pt]
  \item[\textbf{أ)}] يُعفى التلميذ من كل الأنشطة الكتابية نهائيًا، ويُكتفى بتقويمه شفهيًا، لأن الأدوات الصوتية لا تستطيع تحويل كلامه إلى نص، ولا يمكن تكييف الأنشطة بأي طريقة أخرى.
  \item[\textbf{ب)}] الزميل محق، لأن كل استعمال للأداة في التقويم غش، حتى لو كان تكييفًا لصعوبة مثبتة، فيُقوَّم التلميذ بالكتابة اليدوية وحدها مثل زملائه دون أي مراعاة لحالته، تحقيقًا للمساواة الكاملة بين التلاميذ.
  \item[\textbf{ج)}] هو تكييف مشروع يتيح للتلميذ إظهار كفاءته في التعبير، بشرط تحديد الأنشطة المعنية وأهدافها، ومراجعة ما تكتبه الأداة معه، والتنسيق مع الإدارة والأسرة، مع مواصلة دعم مهارة الكتابة اليدوية.
  \item[\textbf{د)}] يُسمح بالإملاء على الأداة، وتُعتمد النصوص التي تكتبها كما هي دون مراجعتها مع التلميذ، لأن الأداة تكتب بالضبط ما ينطقه، فلا تضيف ولا تحذف أي كلمة.
  \item[\textbf{هـ)}] يُسمح له بإملاء كل إجاباته في كل الأنشطة، بما فيها أنشطة الخط والإملاء، دون مراجعة ما تكتبه الأداة، لأن التكييف يعني الإعفاء الكامل من كل ما يتعلق بالكتابة، ويرتاح التلميذ من صعوبته نهائيًا.
\end{itemize}
\par\vspace{6pt}
\begin{center}
{\khatt\textbf{جدول الإجابات}}\\[3pt]
\renewcommand{\arraystretch}{1.6}
\begin{tabular}{|>{\centering\arraybackslash}m{1.9cm}|*{5}{>{\centering\arraybackslash}m{1.2cm}|}}
\hline
\textbf{رقم السؤال} & \textbf{1} & \textbf{2} & \textbf{3} & \textbf{4} & \textbf{5} \\ \hline
\textbf{الإجابة} &  &  &  &  &  \\[8pt] \hline
\end{tabular}

\vspace{8pt}
\renewcommand{\arraystretch}{2.2}
\begin{tabular}{|>{\centering\arraybackslash}m{7.5cm}|>{\centering\arraybackslash}m{3.2cm}|>{\centering\arraybackslash}m{3.2cm}|}
\hline
\textbf{الاسم واللقب} & \textbf{الفوج} & \textbf{العلامة} \\ \hline
 & & \hfill /20 \\ \hline
\end{tabular}
\end{center}
\end{document}
